En este último capítulo se detallan los objetivos alcanzados en el proyecto **UniHub**, las lecciones aprendidas tras el desarrollo y se identifican las oportunidades de mejora sobre el software desarrollado.

\section{Objetivos alcanzados}
Se han alcanzado con éxito el 100\% de los objetivos marcados para el proyecto a lo largo de sus cuatro fases de ingeniería:
\begin{enumerate}
    \item Construcción de un rastreador modular en Python (Fase 1) estructurado en 2 procesos paralelos (Productor de red I/O y Consumidor de análisis CPU) y 2 partes (RUCT/BOE y rastreo de webs oficiales con \texttt{robots.txt}, Sitemap XML y sinónimos), capaz de procesar el catálogo oficial, clasificar Doctorados (RD 99/2011), aplicar filtro estricto de vigencia, escanear todos los PDFs candidatos del BOE, guardar URLs directas en \texttt{web\_fuente\_directa\_url}, aplicar resiliencia de conexión (5m/15m cortocircuito), medir rendimiento y notificar automáticamente a la API REST.
    \item Implementación de un modelo relacional en PostgreSQL (Fase 2) con usuario de solo lectura (\texttt{rupt\_api\_user}), pipeline ETL, ordenación prioritaria (Públicas primero, Privadas después), soporte de métricas físicas de contenedores cgroup y servicio REST en FastAPI con endpoints de consulta y gestión CRUD.
    \item Creación de una aplicación web SPA en React (Fase 3) bajo el sistema de diseño e identidad visual de la Universidad de Cádiz (UCA), paginación configurable (5, 10, 20, 50, 100 resultados por página), ocultación de códigos internos, geolocalización por fórmula de Haversine con distancias integradas en tarjetas, sección de universidades destacadas (top 6 más visitadas), apartado "Sobre Nosotros" (TFG Alejandro Ramos Rodríguez, UCA) y panel de administración aislado en la ruta manual \texttt{/admin}.
    \item Orquestación completa en contenedores Docker independientes (Fase 4) con persistencia permanente de datos relacionales y de archivos en el sistema anfitrión.
\end{enumerate}

\section{Lecciones aprendidas}
Entre las principales lecciones aprendidas destacan:
\begin{itemize}
    \item \textbf{Tolerancia a Fallos y Resiliencia en Scraping Web:} La importancia de aislar excepciones por registro e implementar pausas estratégicas de resiliencia (5m tras 10 fallos y cortocircuito a los 15m) con registro estructurado en \texttt{errores\_crawler.json}.
    \item \textbf{Gestión Estricta de la Integridad de Datos:} La necesidad de filtrar titulaciones extinguidas y descartar valores indeterminados (\texttt{""}, \texttt{"undefined"}) para prevenir la corrupción del repositorio histórico de datos.
    \item \textbf{Diseño Orientado al Usuario y Accesibilidad:} El uso de un sistema de diseño estructurado (colores corporativos UCA, badge rosa para Doctorados) y elementos modulados de paginación para ofrecer una experiencia intuitiva.
    \item \textbf{Arquitectura de Contenerización Aislada:} La configuración adecuada de redes puente virtualizadas y volúmenes nominados para asegurar el desacoplamiento de servicios y la persistencia de datos.
\end{itemize}

\section{Trabajo futuro}
Como líneas de desarrollo futuro se proponen:
\begin{itemize}
    \item Incorporar modelos de lenguaje (LLMs) para el análisis semántico avanzado de competencias en las guías docentes del BOE.
    \item Implementar notificaciones automáticas por correo electrónico cuando una titulación sufra una modificación oficial publicada en el BOE.
    \item Ampliar las analíticas del panel de administración con comparativas interuniversitarias de notas de corte y empleabilidad.
\end{itemize}
